\section{Conclusion}
When aircraft fly close to one another, their exhaust plumes can overlap, leading to the accumulation of emissions in local airspace. In the UTLS, where aircraft typically spend most of their flight, high emissions concentrations of \ce{NO_x} and water vapour can lead to reductions in both ozone production and persistent contrail generation.

This paper initiated an investigation into the chemical response of the atmosphere to aircraft emissions saturation, in high-density airspace. This included a simple kinetic comparison, a coarse resolution analysis on a global scale, and a higher resolution box model study. The calculation using simple kinetic model shows that at cruise altitude, the amount of \ce{NO2} required for parity in OH loss with loss due to reaction with CO and \ce{CH4} is around 2~ppb. Similar calculation using the STOCHEM-CRI global model, and IAGOS measurement data shows that 0.5~ppb of ${[\ce{NO2}]_{\mathrm{equiv}}}$ may exist in northern midlatitudes. Northern midlatitudes would be most likely to benefit from formation flying to suppress ozone production. The box model results show that the impact of aircraft emissions on spatially averaged \ce{O3} could be halved if the aircraft could maintain separations inside 4~km rather than 10~km.

The results generated from these simulations indicate a clear trend toward a climate beneficial outcome due to emissions saturation. In future papers, we aim to increase sophistication of the box model by integrating higher level plume modelling methods and a contrail model into the analysis. In addition to this, the aim is to integrate air traffic trajectories into the study, such that emissions from either real world or simulated flights can be modelled more accurately.